%!TEX root = ../thesis.tex
\mychapter{Conclusion}{Conclusion}{}
\label{chap:conclusion}

\mysection{Answers to the research questions}{Answers to the Research Questions}
\label{sec:answers}

\paragraph{RQ1.} On the physical rover, the local pipeline handled levels 0 to 2
without help in all ten runs, and levels 3 to 5 ended in a completed mission in
11 of 13 runs (nine directly and two after guidance), with one run halted after
guidance at level~3 and one aborted at level~4. Its main weakness at level~4 was
dropping the condition in two of seven plans, and its main weakness at level~3
was speech-to-text noise, which twice changed or lost the meaning of the
command. Reliability at levels 3 to 5 is therefore supported by a handful of
runs and should be reported as a demonstration of capability, not as a rate.
The part of the design that earns its place is the guidance loop, which turned
several runs that would have ended with `target not found' into completed
missions.

\paragraph{RQ2.} In the virtual tests no hallucinated step appeared at any
depth tested, including the extended context sweep run out to 120 chained
instructions. Across all 270 of those runs, the planner never substitutes a
wrong object, location or action. Compositional depth does have a hard limit,
but the failure mode is truncation, not hallucination. As seen from depth 40 onward,
every run at every context setting tested (1024, 2048 and 8192 tokens) stops
generating steps partway through an otherwise-correct plan, and none recovers
at any depth up to 120. Context window size does not the outcome. 
The smaller context settings produce a completely
empty plan in roughly 40\% of deep runs, while the largest setting never does
and sustains close to double the plan length on average. This points to a
generation-length ceiling in the planner's own decoding behaviour rather than
the model running out of context-window capacity to track the room and the
command. Neither this truncation nor the silent dropping of unresolvable steps
is caught by the pipeline at plan time. This argues for a capability and
coverage check on the plan before departure as the most useful next
improvement.

\mysection{Contributions}{Contributions}
\label{sec:contributions}

\mysection{Limitations}{Limitations}
\label{sec:limitations}

\paragraph{Evidence base.} The evidence is one room, one rover, and three to
seven runs per physical level. Level~5 was the least settled of the six. Looking at 
its ten runs, three completed outright, three only completed after a guidance
turn, and four halted or aborted. That is the most mixed result of any level,
so I do not consider it adequately covered yet. The virtual tests has no odometry drift
and use typed text, so they separate plan quality from perception and
speech, and they are not directly comparable to the physical runs. The
ultrasonic stop does not protect against low obstacles approached at an
angle.

\paragraph{Camera and lighting.} The camera, not the planner, was the weak
point of the physical setup. Across levels~3 and~5, every halt with an
operator note attached names poor lighting or low contrast as the reason. These were the following
issues, a black bag invisible against a dark floor, an orange box blending into a
cupboard door, a frame thrown out as too blurry because the room was too dim,
and one run that only worked once the operator physically moved the target
into better light. Level~5 shows this best: the plans themselves were judged
good in every run I reviewed, and the failures were the vision model running
out of usable light once the rover got close enough to need a fine judgement.
This is a limitation of the webcam-and-lighting setup I tested with, not of
the planning pipeline. A better-lit room, or a camera that actually works in
low light, would probably fix most of it.

\paragraph{Reference comparison scope.} The cloud-model comparison
(Section~\ref{sec:cloud_comparison}) only covers the depth-1-to-12 virtual
chains. It was never run on the depth-4-to-120 context sweep, and I never ran
a cloud model on the physical rover. So the finding that a small local model
matches a cloud model on correctness and latency only holds for short virtual
chains. It says nothing about the compositional-depth ceiling, or about real
perception and speech.

\paragraph{Silent plan modification.} The probe chains showed that when a
step can't be carried out, the pipeline's default behaviour is to just drop
it from the plan without saying so: in all 9 cloud runs and 12 of 15 local
runs, not one got flagged to the operator. None of those dropped steps was
recorded as unresolved by the pipeline's own logging either
(\texttt{probe\_logged} came back false every time), so right now the system
has no way to tell a silently shortened plan apart from one that actually
finished. That is as much a gap in my logging as it is in the plan itself,
and it is the clearest concrete thing the capability check proposed under RQ2
above should catch.


\mysection{Recommendations for future work}{Recommendations for Future Work}
\label{sec:future_work}
